% ============================================================
% RÉFÉRENTIEL DES ENVIRONNEMENTS
% ============================================================

\section{Référentiel des environnements}

Ce document présente les différents environnements disponibles dans le projet, avec un exemple d'utilisation pour chacun.

% ------------------------------------------------------------
% Objectifs
% ------------------------------------------------------------

\subsection{Objectifs}

\begin{objectifs}
À la fin de cette partie, vous devez être capable de :
\begin{itemize}
\item comprendre les notions essentielles ;
\item appliquer les résultats du cours ;
\item résoudre les exercices proposés.
\end{itemize}
\end{objectifs}

% ------------------------------------------------------------
% Rappel de cours
% ------------------------------------------------------------

\subsection{Rappel de cours}

\begin{rappelbox}
Un rappel de cours permet de présenter une propriété, une formule ou une notion importante à retenir.

Par exemple, pour tous réels $a$ et $b$ :

(a+b)2=a2+2ab+b2.

\end{rappelbox}

% ------------------------------------------------------------
% Théorème
% ------------------------------------------------------------

\subsection{Théorème}

\begin{theoremebox}
\bb{Théorème de Pythagore.}
Dans un triangle rectangle, le carré de l'hypoténuse est égal à la somme des carrés des deux autres côtés.

AB2+AC2=BC2.

\end{theoremebox}

% ------------------------------------------------------------
% Lemme
% ------------------------------------------------------------

\subsection{Lemme}

\begin{lemmabox}
\bb{Lemme.}
Si $a$ et $b$ sont deux nombres réels tels que $a \leq b$, alors :

a+c≤b+c

pour tout réel $c$.
\end{lemmabox}

% ------------------------------------------------------------
% Définition
% ------------------------------------------------------------

\subsection{Définition}

\begin{definitionbox}
\bb{Définition.}
Une fonction $f$ est dite croissante sur un intervalle $I$ si, pour tous
$x,y \in I$ tels que $x \leq y$, on a :

f(x)≤f(y).

\end{definitionbox}

% ------------------------------------------------------------
% Point méthode
% ------------------------------------------------------------

\subsection{Point méthode}

\begin{methodebox}
Pour résoudre une équation du premier degré :
\begin{enumerate}
\item développer et réduire les deux membres ;
\item regrouper les termes contenant l'inconnue d'un même côté ;
\item isoler l'inconnue ;
\item vérifier éventuellement la solution obtenue.
\end{enumerate}
\end{methodebox}

% ------------------------------------------------------------
% Attention / erreurs fréquentes
% ------------------------------------------------------------

\subsection{Attention}

\begin{attentionbox}
Attention aux changements de signe lors du passage d'un terme d'un membre à l'autre.

Par exemple :

x+3=7

donne

x=7−3=4,

et non $x=7+3$.
\end{attentionbox}

% ------------------------------------------------------------
% Exercices
% ------------------------------------------------------------

\subsection{Exercice}

\begin{exercicebox}
Résoudre dans $\mathbb{R}$ l'équation :

3x+5=17.

\end{exercicebox}

% ------------------------------------------------------------
% Correction
% ------------------------------------------------------------

\subsection{Correction}

\begin{correctionbox}
On part de :

3x+5=17.

On soustrait $5$ aux deux membres :

3x=12.

On divise par $3$ :

x=4.

La solution est donc :

x=4.

\end{correctionbox}
